\documentclass[11pt,a4paper]{article}
\usepackage[utf8]{inputenc}
\usepackage[T1]{fontenc}
\usepackage[french]{babel}
\usepackage{lmodern}
\usepackage{geometry}
\geometry{margin=2.5cm}
\usepackage{titlesec}
\usepackage{xcolor}
\usepackage{enumitem}
\usepackage{hyperref}

% Personnalisation des titres
\titleformat{\section}{\Large\bfseries\color{blue!80!black}}{\thesection}{1em}{}
\titleformat{\subsection}{\large\bfseries\color{blue!60!black}}{\thesubsection}{1em}{}

\title{\textbf{Argumentaire d'Innovation :\\ La Sécurité par la Data \& l'Intelligence Artificielle\\ dans la Logistique du Dernier Kilomètre}}
\author{Ingénierie Data \& Intelligence Artificielle}
\date{\today}

\begin{document}

\maketitle

\section*{Note de réalité (état livré)}

Ce document justifie l'axe retenu ; il a été rédigé avant le développement et
décrit l'intention. L'état réellement livré et vérifiable est détaillé dans
\texttt{argumentaire\_securite\_realiste.tex}. Trois écarts, assumés :

\begin{itemize}
    \item \textbf{HMAC, et non SHA-256} : un simple hash, recalculable par les
    deux parties, ne prouve pas l'anti-répudiation. Le scellement livré est un
    HMAC-SHA256 sous une clé détenue par le serveur, dérivée du secret racine et
    distincte de la clé de signature des JWT.
    \item \textbf{Évaluation dans la JVM, et non un microservice Python} : le
    modèle de forêt aléatoire est entraîné hors-ligne en Python
    (\texttt{ml/train\_fraud\_model.py}) puis évalué dans le processus Spring Boot
    au fil des positions GPS. Un seul artefact à déployer, au lieu de deux.
    \item \textbf{Row-Level Security non implémenté} : le cloisonnement par RLS
    suppose une refonte du modèle de connexion (pool JPA mono-utilisateur) ; il
    est listé comme hors périmètre dans l'argumentaire vérifiable. L'isolation
    par rôle est assurée au niveau applicatif et contrôlée par des tests.
\end{itemize}

\section*{Introduction}
Dans le cadre de la conception d'une plateforme de livraison (SwiftDeliver), se contenter des standards de sécurité classiques (Authentification JWT, HTTPS, hachage des mots de passe) ne constitue plus une innovation, mais un prérequis. Pour un ingénieur spécialisé en \textbf{Data Engineering et Intelligence Artificielle}, la véritable innovation réside dans le croisement entre la cybersécurité et la science des données : la \textbf{Cybersécurité Comportementale et Spatiale}. 

Ce document justifie pourquoi l'axe de la sécurité, lorsqu'il est traité sous l'angle de la Data et de l'IA, est le domaine le plus pertinent et le plus innovant à explorer pour ce projet.

\section{La Sécurité Moderne est un Problème de Data}
Dans les plateformes logistiques (Uber, Deliveroo), les failles de sécurité majeures ne proviennent pas d'attaques informatiques traditionnelles (injections SQL), mais de la \textbf{fraude comportementale} : utilisation d'applications "Fake GPS" par les livreurs, vols de marchandises déguisés en retards, ou falsification de preuves de livraison. 
Résoudre ces failles nécessite l'ingestion et l'analyse de flux massifs de données télémétriques (IoT). Ainsi, le défi sécuritaire devient un défi pur de \textbf{Data Engineering} (Stream Processing) et de \textbf{Machine Learning} (Anomaly Detection).

\section{Axes d'Innovation Sécuritaire par la Data et l'IA}

\subsection{A. Détection d'Anomalies Spatiales (Anti-GPS Spoofing)}
Les téléphones des livreurs agissent comme des capteurs IoT générant un flux de données en temps réel.
\begin{itemize}
    \item \textbf{L'approche classique :} Sauvegarder passivement la position.
    \item \textbf{Retenue :} comparer en temps réel le flux GPS brut au flux attendu vers la destination (graphe routier OSRM). Une forêt aléatoire, entraînée hors-ligne et évaluée dans la JVM au fil des positions (voir la note de réalité), analyse la physique du mouvement pour détecter les "téléportations" (GPS Spoofing), les déviations injustifiées (suspicion de vol) ou les arrêts prolongés. L'IA protège ainsi l'intégrité physique et numérique des opérations.
\end{itemize}

\subsection{B. Preuve de Livraison Cryptographique (Proof of Delivery)}
La répudiation (le client affirme ne rien avoir reçu, le livreur affirme avoir livré) est une vulnérabilité métier critique.
\begin{itemize}
    \item \textbf{Retenu, sous forme renforcée :} à l'instant de la livraison, le système scelle par HMAC-SHA256 (et non un simple hash recalculable) les coordonnées GPS du livreur, celles du client, et un horodatage précis. Une validation algorithmique bloque la transaction si la distance excède un seuil de sécurité. C'est l'application de la cryptographie aux données spatiales pour garantir l'anti-répudiation.
\end{itemize}

\subsection{C. Zero Trust et Isolation des Données (Row-Level Security)}
En complément de l'IA, l'architecture des données doit être nativement sécurisée.
\begin{itemize}
    \item \textbf{Piste non retenue (hors périmètre) :} implémenter le \textit{Row-Level Security} (RLS) directement dans PostgreSQL. Séduisant sur le principe, il suppose une refonte du modèle de connexion (pool JPA mono-utilisateur) et n'a pas été retenu ; l'isolation est assurée au niveau applicatif (contrôle par rôle et par propriété), contrôlée par des tests.
\end{itemize}

\section{Adéquation avec le Profil Data Engineering \& IA}
Se concentrer sur cet aspect sécuritaire permet de démontrer l'ensemble du spectre de compétences attendu :
\begin{enumerate}
    \item \textbf{Architecture de flux (Data Engineering) :} Captation, traitement en temps réel (Redis) et stockage de données IoT massives.
    \item \textbf{Intelligence Artificielle (Machine Learning) :} Modélisation de comportements normaux et détection d'anomalies (Outlier Detection) sur des séries temporelles spatiales.
    \item \textbf{Systèmes Distribués (Cloud) :} Déploiement de microservices capables de soutenir l'analyse sécuritaire sans impacter les performances de l'application.
\end{enumerate}

\section{Planification des Sprints (Implémentation Locale \& Gratuite)}
L'intégralité de cette architecture a été pensée pour être \textbf{100\% gratuite et exécutable localement}, sans nécessité de déploiement Cloud payant pour la soutenance (grâce à l'utilisation de Docker, d'Ollama pour l'IA, et d'OSRM hébergé localement). Voici le plan d'implémentation :

\begin{itemize}
    \item \textbf{Sprint 1 : Preuve de Livraison Cryptographique (PoD) - [Semaine 1]}
    \begin{itemize}
        \item Génération d'un hash SHA-256 dans le backend Spring Boot lors du clic "Livré".
        \item Calcul de la distance entre le livreur et la destination pour autoriser/bloquer la transaction.
    \end{itemize}
    \item \textbf{Sprint 2 : Pipeline IoT \& Détection d'Anomalies Spatiales - [Semaine 2-3]}
    \begin{itemize}
        \item Création du microservice Python (Data Worker).
        \item Ingestion du flux GPS depuis Redis (Pub/Sub).
        \item Comparaison algorithmique des trajectoires réelles vs la \textit{polyline} optimale OSRM locale (détection des détours et arrêts).
    \end{itemize}
    \item \textbf{Sprint 3 : Sécurisation Zero-Trust de la Base de Données - [Semaine 4]}
    \begin{itemize}
        \item Activation de \textit{Row-Level Security} (RLS) sur PostgreSQL.
        \item Création des politiques de sécurité pour scinder hermétiquement les données entre clients, livreurs et administrateurs.
    \end{itemize}
\end{itemize}

\section*{Conclusion}
Innover dans la sécurité ne signifie pas faire de la cryptographie fondamentale, mais utiliser \textbf{l'intelligence des données pour protéger l'écosystème}. En choisissant la détection de fraude par l'IA et l'analyse de flux spatiaux, le projet répond à un problème critique de cybersécurité logistique tout en exploitant au maximum les compétences en Data Engineering et en Intelligence Artificielle. C'est le meilleur compromis pour démontrer une expertise technique transversale de haut niveau.

\end{document}
